%! Author = sriadityavedantam
%! Date = 4/8/26

\section{Rational Differential 1-Forms}

Let $X$ be a quasi-projective variety.
A function $f \in \c(X)$ is rational if it is regular on some open dense subset $U \subseteq X$, i.e.\ $f \in \c[U]$.
Analogously, define $\Omega(X)$ to be the space of all rational differential $1$-forms on $X$.
The choice of $U$ does not matter:
if $\omega \in \Omega(U)$ and $\omega' \in \Omega(U')$ agree on $U \cap U'$, then they define the same rational $1$-form on $X$.
The space $\Omega(X)$ is a $\c(X)$-vector space.
Indeed, for $\omega, \omega' \in \Omega(X)$ and $f \in \c(X)$,
\[
    \omega + \omega' \in \Omega(X),
    \quad
    f\omega \in \Omega(X).
\]

\subsection*{Dimension}

We claim
\[
    \dim_{\c(X)} \Omega(X) = \dim(X).
\]

\subsection*{Affine Case}

Let $X = \aa^n$ with coordinates $t_1,\ldots,t_n$.
A rational $1$-form has the form
\[
    \sum_i f_i \, d g_i,
    \quad
    f_i, g_i \in \c(X).
\]

Then
\[
    \{dt_1,\ldots,dt_n\}
\]
forms a basis for $\Omega(X)$ over $\c(X)$.
Hence
\[
    \dim_{\c(X)} \Omega(X) = n = \dim(X).
\]

\subsection*{Projective Case}

If $X = \pp^n$, then
\[
    \Omega[\pp^n] = 0,
\]
since all regular functions are constant.
However,
\[
    \Omega(\pp^n)
\]
is an $n$-dimensional vector space over $\c(\pp^n)$.
Moreover,
\[
    \Omega(\pp^n) \cong \Omega(\aa^n)
\]
on affine charts.

\subsection*{Birational Invariance}

In general, if $X$ and $X'$ are birational varieties, then there is a natural identification
\[
    \Omega(X) \cong \Omega(X').
\]